\documentclass[11pt]{article}

\usepackage{amsmath, amssymb, amsthm}
\usepackage{geometry}
\usepackage{graphicx}
\usepackage{booktabs}

\geometry{margin=1in}

\title{CSE400: Fundamentals of Probability in Computing\\
Lecture Scribe (L15--L18)\\
Joint Random Variables and Joint Distributions}
\author{Name: Masoom Choksi \\ AUID: AU2440263}
\date{March 10, 2026}

\begin{document}

\maketitle

\section{Outline}

\begin{itemize}
\item Why Study Joint Random Variables
\item Discrete Case
\begin{itemize}
\item Joint PMF Table Representation
\item Example: Rolling Two Dice
\end{itemize}
\item Continuous Case
\begin{itemize}
\item Geometric Interpretation
\item Worked Example
\end{itemize}
\item Joint Cumulative Distribution Function
\begin{itemize}
\item Definition of Joint CDF
\item Continuous Case
\item Discrete Case
\item Properties of Joint CDF
\end{itemize}
\item Comprehensive Example
\end{itemize}

\section{Why Study Joint Random Variables}

In many real world problems we are interested in studying more than one random variable simultaneously.

Examples include:

\subsection{Smart Transportation}

Let

\[
X = \text{Vehicle Speed}
\]

\[
Y = \text{Traffic Density}
\]

Question:

\[
\text{Can speed be high when traffic density is also high?}
\]

These variables may be correlated.

Applications include:

\begin{itemize}
\item Adaptive traffic signals
\item Congestion prediction
\item Autonomous vehicle routing
\end{itemize}

\subsection{Wireless Network Performance}

Let

\[
X = \text{Signal-to-Noise Ratio (SNR)}
\]

\[
Y = \text{Data Throughput}
\]

Throughput depends on SNR but may also depend on congestion and interference.

Applications include:

\begin{itemize}
\item 5G scheduling
\item Adaptive modulation
\item Network planning
\end{itemize}

\subsection{Additional Examples}

\begin{itemize}
\item Weather Prediction: \(X =\) Rainfall, \(Y =\) Temperature
\item Drone Delivery: \(X =\) Wind Speed, \(Y =\) Delivery Time
\item Rolling Two Dice: \(X\) and \(Y\) denote outcomes of two dice
\end{itemize}

These scenarios require the study of joint random variables.

\section{Discrete Case: Joint Probability Mass Function}

\subsection{Definition}

Let \(X\) and \(Y\) be discrete random variables.

The \textbf{joint probability mass function (joint PMF)} is defined as

\[
p_{X,Y}(x,y) = P(X = x, Y = y)
\]

for all possible values of \(x\) and \(y\).

\subsection{Properties}

\begin{enumerate}
\item Non-negativity:

\[
p_{X,Y}(x,y) \ge 0
\]

\item Normalization:

\[
\sum_x \sum_y p_{X,Y}(x,y) = 1
\]
\end{enumerate}

\subsection{Joint PMF Table Representation}

The joint PMF of two discrete variables can be represented using a table.

\begin{center}
\begin{tabular}{c|ccc}
 & $y_1$ & $y_2$ & $y_3$ \\
\hline
$x_1$ & $p(x_1,y_1)$ & $p(x_1,y_2)$ & $p(x_1,y_3)$ \\
$x_2$ & $p(x_2,y_1)$ & $p(x_2,y_2)$ & $p(x_2,y_3)$ \\
$x_3$ & $p(x_3,y_1)$ & $p(x_3,y_2)$ & $p(x_3,y_3)$
\end{tabular}
\end{center}

Each cell represents the probability of the joint event.

\section{Example: Rolling Two Dice}

Let

\[
X = \text{Outcome of first die}
\]

\[
Y = \text{Outcome of second die}
\]

Each die takes values

\[
\{1,2,3,4,5,6\}
\]

Since the dice are fair:

\[
P(X=x,Y=y) = \frac{1}{36}
\]

for all

\[
x,y \in \{1,2,3,4,5,6\}
\]

\subsection{Example Event}

Compute probability that

\[
X + Y = 7
\]

Possible outcomes:

\[
(1,6), (2,5), (3,4), (4,3), (5,2), (6,1)
\]

Number of outcomes:

\[
6
\]

Thus

\[
P(X+Y=7) = \frac{6}{36} = \frac{1}{6}
\]

\section{Continuous Case: Joint Probability Density Function}

\subsection{Definition}

Let \(X\) and \(Y\) be continuous random variables.

The \textbf{joint probability density function (joint PDF)} is denoted

\[
f_{X,Y}(x,y)
\]

such that

\[
P((X,Y) \in A) =
\int\int_A f_{X,Y}(x,y)\,dx\,dy
\]

\subsection{Properties}

\begin{enumerate}
\item Non-negativity:

\[
f_{X,Y}(x,y) \ge 0
\]

\item Normalization:

\[
\int_{-\infty}^{\infty}\int_{-\infty}^{\infty} f_{X,Y}(x,y)\,dx\,dy = 1
\]
\end{enumerate}

\subsection{Geometric Interpretation}

For continuous variables, probability corresponds to the volume under the joint density surface.

For a region \(A\):

\[
P((X,Y) \in A) =
\int\int_A f_{X,Y}(x,y)\,dx\,dy
\]

\section{Worked Example}

Suppose

\[
f_{X,Y}(x,y) =
\begin{cases}
c & 0 \le x \le 1,\ 0 \le y \le 1 \\
0 & \text{otherwise}
\end{cases}
\]

\subsection{Step 1: Determine Constant \(c\)}

Using normalization:

\[
\int_0^1 \int_0^1 c \, dx\,dy = 1
\]

\[
c \int_0^1 \int_0^1 dx\,dy = 1
\]

\[
c (1)(1) = 1
\]

\[
c = 1
\]

\subsection{Step 2: Probability Calculation}

Compute

\[
P(X < 0.5, Y < 0.5)
\]

\[
P =
\int_0^{0.5} \int_0^{0.5} 1 \, dx\,dy
\]

\[
P = (0.5)(0.5)
\]

\[
P = 0.25
\]

\section{Joint Cumulative Distribution Function}

\subsection{Definition}

The \textbf{joint cumulative distribution function (Joint CDF)} of random variables \(X\) and \(Y\) is defined as

\[
F_{X,Y}(x,y) = P(X \le x, Y \le y)
\]

\section{Joint CDF for Continuous Variables}

If the joint PDF exists:

\[
F_{X,Y}(x,y) =
\int_{-\infty}^{x}
\int_{-\infty}^{y}
f_{X,Y}(u,v)\,dv\,du
\]

\section{Joint CDF for Discrete Variables}

For discrete variables:

\[
F_{X,Y}(x,y) =
\sum_{u \le x}
\sum_{v \le y}
p_{X,Y}(u,v)
\]

\section{Properties of Joint CDF}

\begin{enumerate}
\item Range:

\[
0 \le F_{X,Y}(x,y) \le 1
\]

\item Non-decreasing in both arguments.

\item Limits:

\[
\lim_{x \to -\infty} F_{X,Y}(x,y) = 0
\]

\[
\lim_{y \to -\infty} F_{X,Y}(x,y) = 0
\]

\[
\lim_{x \to \infty, y \to \infty} F_{X,Y}(x,y) = 1
\]
\end{enumerate}

\section{Summary}

\begin{itemize}
\item Joint random variables describe the probabilistic relationship between two variables.
\item For discrete variables, the relationship is captured by the joint PMF.
\item For continuous variables, the relationship is captured by the joint PDF.
\item Probabilities correspond to summations (discrete) or double integrals (continuous).
\item The joint CDF provides the cumulative probability of both variables simultaneously.
\end{itemize}

\end{document}